\label{sec:related_work}

A comprehensive survey of proposed methods for different action recognition tasks can be found in \cite{HEO202632, KongFu2022} and a nicely concise history of development is included in \cite{shuchang2022surveyhumanactionrecognition}.
Hence, we shall confine ourselves to some recent relevant models herein.

\paragraph{SlowFast Networks \cite{feichtenhofer2019slowfast}.}
This network adopts the Two–stream CNN–based framework \cite{simonyan2014twostream} for incorporation between spatial and temporal semantics of video data.
It was originally inspired by the biological studies in primate visual system, and prematurely mimicking the mechanisms of Parvocellular and Magnocellular to develop the Slow-Fast pathways and lateral connections between them.
Evidently, different variants achieved, on Kinetic–400 \cite{kay2017kinetics}: 75.6\% to 79.8\% top–1 accuracy, 92.1\% to 93.9\% top–5 accuracy, on Kinetic–600 \cite{carreira2018shortnotekinetics600}: 78.8\% to 81.8\%, and on Charades \cite{Charades}: 39.0 to 45.2 mAP.

\paragraph{EfficientNet \cite{tan2019efficientnet}.} This model relies on the scalability philosophy of DL models, e.g. Residual Networks \cite{ResNet}, and platform–aware search of network architecture \cite{MnasNet}.
Then, by scaling the ConvNets wisely across three dimensions, i.e. depth, width and resolution. Different model variants achieved, on ImageNet dataset \cite{Russakovsky2015}, 77.1\% to 84.3\% top–1 accuracy and 93.3\% to 97.0\% top–5 accuracy, with fewer learnable parameters, from 5.3M to 66M.
The \textbf{X3D} model was built upon the philosophy in this network type, but offers more computational advantages.

\paragraph{MViT \cite{fan2021mvit,li2022mvitv2}.} This model was established by leveraging the \textit{multiscale feature hierarchies} in CNN–based models and the \textit{attention} mechanism in Transformer–based models.
By doing this, the model can establish a pyramid of features similar to CNN while globalising the interactions amongst the local regions thanks to self–attention.
Therefore, it reduces the locality bias in CNN–based models and simultaneously breaking the permutation invariance in ViT \cite{ViT}, which is semantically wrong in the context of video understanding where temporal relationship must be considered.
As a result, it gained, on Kinetic–400 \cite{kay2017kinetics}, 76.0\% to 81.2\% top–1 accuracy and 92.1\% to 95.1\% top–5 accuracy, on Kinetic–600 \cite{carreira2018shortnotekinetics600}, 82.1\% to 83.8\% top–1 accuracy and 95.7\% to 96.3\% top–5 accuracy,
and on Charades \cite{Charades}, 43.9 to 47.7 mAP.

\paragraph{MoViNets \cite{MoViNets}.} This model substantially reduces the peak memory usage and computational budget while performing competitively in the benchmark datasets above.
The authors utilise the neural architecture search to determine a good construction of 3D CNN architectures. Then, they introduce stream buffer component to simultaneously reduce memory consumption and maintain long-range temporal dependencies of sub–clips.
Still, it loses some accuracy, by roughly 1\% in Kinetic-600 \cite{carreira2018shortnotekinetics600}. Thus, they recover the accuracy loss by leveraging the ensembling philosophy \cite{kondratyuk2020ensembling}, specifically, they train two model variants independently and combine them by taking the arithmetic mean on the unweighted logits before their final softmax layer.

Indeed, some models, in the family of MoViNets, are better, in terms of GFLOPs and top-1 accuracy, than certain variants of X3D, as reported in Appendix C of \cite{MoViNets}.
Nonetheless, it was originally implemented in TensorFlow, not PyTorchVideo, though there is an incomplete unofficial implementation, can be found in \href{https://github.com/Atze00/MoViNet-pytorch.git}{\texttt{this link}}.
We shall confine ourselves to \textbf{X3D} due to its complete availability and technical ease.


